\textbf{Phần 1. Trắc nghiệm lựa chọn}
\setcounterpageref{ex}{0}

\begin{ex} %[3D1V76]
Số liền trước của số bé nhất có năm chữ số là
\choice
{$10\,001$}
{$10\,000$}
{$9990$}
{\True $9999$}
\loigiai{
Số bé nhất có 5 chữ số là $10\,000$. Số liền trước là: $10\,000 - 1 = 9999$.
\textbf{Đáp án của câu hỏi là D}
}
\end{ex}

\begin{ex} %[3D1V76]
Số gồm 8 chục nghìn, 5 trăm và 4 đơn vị được viết là
\choice
{$85\,004$}
{\True $80\,504$}
{$80\,540$}
{$85\,400$}
\loigiai{
Hàng nghìn và hàng chục là 0: $80\,504$.
\textbf{Đáp án của câu hỏi là B}
}
\end{ex}

\begin{ex} %[3D1V76]
Làm tròn số $64\,850$ đến hàng nghìn ta được số
\choice
{$64\,000$}
{$64\,900$}
{\True $65\,000$}
{$60\,000$}
\loigiai{
Hàng trăm là 8 nên làm tròn lên $65\,000$.
\textbf{Đáp án của câu hỏi là C}
}
\end{ex}

\begin{ex} %[3D1V76]
Số La Mã $\text{XIX}$ biểu thị số tự nhiên nào?
\choice
{\True $19$}
{$21$}
{$18$}
{$14$}
\loigiai{
$\text{XIX} = 19$.
\textbf{Đáp án của câu hỏi là A}
}
\end{ex}

\begin{ex} %[3D1V76]
Số chẵn bé nhất có năm chữ số khác nhau là
\choice
{$10\,246$}
{$10\,000$}
{$12\,346$}
{\True $10\,234$}
\loigiai{
Số bé nhất có 5 chữ số khác nhau là $10\,234$ (tận cùng 4 là số chẵn).
\textbf{Đáp án của câu hỏi là D}
}
\end{ex}

\begin{ex} %[3D1V76]
Số $100\,000$ gấp số $1000$ bao nhiêu lần?
\choice
{$10$ lần}
{\True $100$ lần}
{$1000$ lần}
{$10\,000$ lần}
\loigiai{
$100\,000 : 1000 = 100$ (lần).
\textbf{Đáp án của câu hỏi là B}
}
\end{ex}

\begin{ex} %[3D1V76]
Từ năm chữ số $0; 1; 3; 6; 8$, viết được số có năm chữ số khác nhau lớn nhất là
\choice
{$86\,301$}
{$86\,130$}
{\True $86\,310$}
{$83\,610$}
\loigiai{
Thứ tự giảm dần: $86\,310$.
\textbf{Đáp án của câu hỏi là C}
}
\end{ex}

\begin{ex} %[3D1V76]
Số liền sau của số lớn nhất có năm chữ số là
\choice
{\True $100\,000$}
{$99\,998$}
{$100\,001$}
{$99\,990$}
\loigiai{
Số lớn nhất có 5 chữ số là 99 999. Số liền sau là: $99\,999 + 1 = 100\,000$.
\textbf{Đáp án của câu hỏi là A}
}
\end{ex}

\begin{ex} %[3D1V76]
Làm tròn số $75\,000$ đến hàng chục nghìn ta được số
\choice
{$70\,000$}
{$75\,000$}
{$90\,000$}
{\True $80\,000$}
\loigiai{
Hàng nghìn là 5 nên làm tròn lên $80\,000$.
\textbf{Đáp án của câu hỏi là D}
}
\end{ex}

\begin{ex} %[3D1V76]
Trong số $78\,942$, chữ số 7 có giá trị là
\choice
{$7000$}
{\True $70\,000$}
{$700$}
{$700\,000$}
\loigiai{
Chữ số 7 ở hàng chục nghìn có giá trị là $70\,000$.
\textbf{Đáp án của câu hỏi là B}
}
\end{ex}

\begin{ex} %[3D1V76]
Số nào sau đây bé hơn $54\,320$?
\choice
{$54\,325$}
{$54\,400$}
{\True $54\,302$}
{$55\,000$}
\loigiai{
$54\,302 < 54\,320$.
\textbf{Đáp án của câu hỏi là C}
}
\end{ex}

\begin{ex} %[3D1V76]
Dãy số nào sau đây được xếp theo thứ tự giảm dần?
\choice
{\True $80\,000; 75\,000; 72\,000; 70\,000$}
{$70\,000; 72\,000; 75\,000; 80\,000$}
{$80\,000; 72\,000; 75\,000; 70\,000$}
{$75\,000; 80\,000; 72\,000; 70\,000$}
\loigiai{
Thứ tự giảm dần chuẩn: $80\,000 > 75\,000 > 72\,000 > 70\,000$.
\textbf{Đáp án của câu hỏi là A}
}
\end{ex}

\textbf{Phần 2. Trắc nghiệm đúng sai}
\setcounterpageref{ex}{0}

\begin{ex} %[3D1V76]
Xét tính đúng sai của các nhận định ôn tập số học:
\choiceTF
{\True Số 100 000 gồm 10 chục nghìn.}
{Chữ số La Mã XXI có giá trị là 19.}
{\True Khi làm tròn số 62 480 đến hàng nghìn thì được 62 000.}
{Số 99 999 là số chẵn.}
\loigiai{
\begin{itemize}
    \item Ý a) Đúng: Cấu tạo số chuẩn.
    \item Ý b) Sai: Sai vì XXI là 21.
    \item Ý c) Đúng: Vì hàng trăm là 4.
    \item Ý d) Sai: Sai vì tận cùng là 9 nên là số lẻ.
\end{itemize}
\begin{itemize}
    \item a) \textbf{Đúng}.
    \item b) \textbf{Sai}.
    \item c) \textbf{Đúng}.
    \item d) \textbf{Sai}.
\end{itemize}
}
\end{ex}

\begin{ex} %[3D1V76]
Cho số $85\,460$:
\choiceTF
{\True Làm tròn đến hàng trăm được số 85 500.}
{\True Làm tròn đến hàng nghìn được số 85 000.}
{\True Làm tròn đến hàng chục nghìn được số 90 000.}
{Số liền sau của số đó là 85 459.}
\loigiai{
\begin{itemize}
    \item Ý a) Đúng: Đơn vị là 6 làm tròn lên.
    \item Ý b) Đúng: Trăm là 4 làm tròn xuống.
    \item Ý c) Đúng: Nghìn là 5 làm tròn lên.
    \item Ý d) Sai: Sai vì số liền sau là 85 461.
\end{itemize}
\begin{itemize}
    \item a) \textbf{Đúng}.
    \item b) \textbf{Đúng}.
    \item c) \textbf{Đúng}.
    \item d) \textbf{Sai}.
\end{itemize}
}
\end{ex}

\textbf{Phần 3. Trả lời ngắn}
\setcounterpageref{ex}{0}

\begin{ex} %[3D1V76]
Tìm số liền sau của số 99 999.
\shortans{100000}
\loigiai{
\textbf{Đáp số:} $100000$\\
$99\,999 + 1 = 100\,000$.
}
\end{ex}

\begin{ex} %[3D1V76]
Viết số tự nhiên tương ứng với chữ số La Mã XVIII.
\shortans{18}
\loigiai{
\textbf{Đáp số:} $18$\\
XVIII là 18.
}
\end{ex}

\begin{ex} %[3D1V76]
Làm tròn số 76 540 đến hàng nghìn.
\shortans{77000}
\loigiai{
\textbf{Đáp số:} $77000$\\
Hàng trăm là 5 nên làm tròn thành 77 000.
}
\end{ex}

\begin{ex} %[3D1V76]
Viết số lớn nhất có năm chữ số khác nhau.
\shortans{98765}
\loigiai{
\textbf{Đáp số:} $98765$\\
Số đó là 98 765.
}
\end{ex}

\begin{ex} %[3D1V76]
Tính tổng các chữ số của số 45 678.
\shortans{30}
\loigiai{
\textbf{Đáp số:} $30$\\
$4 + 5 + 6 + 7 + 8 = 30$.
}
\end{ex}

\begin{ex} %[3D1V76]
Có bao nhiêu số tròn nghìn từ 30 000 đến 40 000?
\shortans{11}
\loigiai{
\textbf{Đáp số:} $11$\\
$(40\,000 - 30\,000) : 1000 + 1 = 11$ số.
}
\end{ex}
